%-------------------------------------- COMIENZA descripción del caso de uso.
\begin{UseCase}{CU-02}{Gestión de usuarios y roles}{
	Permite al Director registrar a los miembros del personal multidisciplinario, asignarles su rol respectivo, actualizar sus datos personales o laborales, y habilitar o deshabilitar sus cuentas de acceso en el sistema.
}
	\UCitem{Versión}{\color{Gray}1.0}
	\UCitem{Autor}{\color{Gray}Ramírez Cuevas Emiliano}
	\UCitem{Supervisa}{\color{Gray}Guerra Salinas Edgar Rafael}
	\UCitem{Actor}{\hyperlink{tDirector}{Director}}
	\UCitem{Propósito}{Administrar las cuentas y roles del personal operativo en el sistema Taimotla para asegurar el correcto flujo de trabajo.}
	\UCitem{Entradas}{CURP, RFC, Nombre, Apellidos, Fecha de nacimiento, Sexo, Correo, Teléfono, Domicilio, Rol, Cédula profesional, Especialidad y Contraseña del nuevo usuario.}
	\UCitem{Origen}{Teclado}
	\UCitem{Salidas}{Confirmación del registro o modificación del usuario en el panel del Director.}
	\UCitem{Destino}{Pantalla y base de datos relacional.}
	\UCitem{Precondiciones}{El Director debe haber iniciado sesión con credenciales válidas en la plataforma.}
	\UCitem{Postcondiciones}{El personal modificado o agregado queda registrado y disponible para la formación de equipos multidisciplinarios.}
	\UCitem{Errores}{
		\begin{description}
			\item[E1:] Datos obligatorios incompletos o en formato incorrecto (ej. CURP no válida).
			\item[E2:] Correo electrónico o CURP duplicados en el sistema.
		\end{description}
	}
 	\UCitem{Tipo}{Caso de uso primario}
 	\UCitem{Observaciones}{Regido por las reglas de negocio \hyperlink{BR-01}{BR-01}, \hyperlink{BR-06}{BR-06}, \hyperlink{BR-07}{BR-07} y \hyperlink{BR-08}{BR-08}.}
 \end{UseCase}
 
 \begin{UCtrayectoria}
 	\UCpaso[\UCactor] Accede al panel de \hyperlink{IU24}{administración de personal de Taimotla}.
 	\UCpaso Despliega la lista de empleados registrados (activos e inactivos).
 	\UCpaso[\UCactor] Presiona el botón \IUbutton{Registrar nuevo personal} o selecciona un empleado para editar.
 	\UCpaso Despliega el formulario de registro/edición de datos de \hyperlink{Personal}{personal}.
 	\UCpaso[\UCactor] Llena los campos correspondientes a la información del personal y presiona el botón \IUbutton{Guardar}.
 	\UCpaso Valida que se hayan capturado todos los campos obligatorios del empleado de acuerdo a su rol, incluyendo la cédula profesional obligatoria para especialistas según la regla de negocio \hyperlink{BR-08}{BR-08}. \Trayref{A}.
 	\UCpaso Verifica que el correo, RFC y CURP ingresados no se encuentren duplicados en la base de datos, y que cumplan con la estructura requerida según las reglas de negocio \hyperlink{BR-06}{BR-06} y \hyperlink{BR-07}{BR-07}. \Trayref{B}.
 	\UCpaso Encripta la contraseña de acceso por seguridad.
 	\UCpaso Guarda la información en la base de datos asociando al empleado con su rol de especialidad respectivo (Abogado, Médico, Psicólogo, Trabajador Social).
 	\UCpaso Despliega el mensaje de éxito \hyperlink{mMSG09}{\bf MSG-09} indicando el registro del usuario.
 	\UCpaso[] -- -- Fin del caso de uso.
 \end{UCtrayectoria}

\begin{UCtrayectoriaA}{A}{Faltan datos de entrada obligatorios o son incorrectos}
	\UCpaso Muestra el mensaje de error \hyperlink{mMSG04}{\bf MSG-04} indicando los campos vacíos o con formato inválido.
	\UCpaso[\UCactor] Oprime el botón \IUbutton{Aceptar}.
	\UCpaso Regresa al paso 5 de la trayectoria principal.
\end{UCtrayectoriaA}

\begin{UCtrayectoriaA}{B}{Duplicidad de datos de identidad en el sistema}
	\UCpaso Muestra el mensaje de error \hyperlink{mMSG02}{\bf MSG-02} alertando que la CURP o correo ya existen en otra cuenta registrada.
	\UCpaso[\UCactor] Oprime el botón \IUbutton{Aceptar}.
	\UCpaso Regresa al paso 5 de la trayectoria principal.
\end{UCtrayectoriaA}
%-------------------------------------- TERMINA descripción del caso de uso.
